\documentclass[11pt]{article}
\usepackage[T1]{fontenc}
\usepackage[utf8]{inputenc}
\usepackage{lmodern}
\usepackage[margin=1in]{geometry}
\usepackage{amsmath,amssymb,amsthm}
\usepackage[hidelinks]{hyperref}
\newtheorem{theorem}{Theorem}
\newtheorem{lemma}[theorem]{Lemma}
\newtheorem{proposition}[theorem]{Proposition}
\theoremstyle{remark}
\newtheorem{remark}[theorem]{Remark}
\newcommand{\Z}{\mathbb Z}
\newcommand{\ind}{\mathbf 1}
\newcommand{\packd}{d_{\mathrm p}}
\newcommand{\ud}{\overline d}

\title{The Packing Density of Five Integers\protect\\with Gaps \(a,a,b,b\)}
\author{}
\date{October 13, 2026}

\begin{document}
\maketitle

\begin{abstract}
For coprime positive integers \(a,b\), we determine the packing density of
\(\{0,a,2a,2a+b,2a+2b\}\). The density is defined using the symmetric
upper density of translation centers, with no periodicity assumption.
The resulting formula is the \(k=m=2\) case of the conjecture following
Theorem~2 of Natalchenko and Sagdeev. An average over six five-point paths
gives a local counting inequality. Equality forces a parity condition,
which supplies the rounding needed in two of the four nonzero residue
classes. A bounded-length interval decomposition then proves the upper
bound for every center set. The matching periodic constructions are due
to Natalchenko and Sagdeev; their specialization is verified here in full.
\end{abstract}

\section{Statement and relation to the source problem}

For integers \(u\le v\), write
\[
 [u,v)=\{z\in\Z:u\le z<v\},\qquad
 [u,v]=\{z\in\Z:u\le z\le v\}.
\]
For a finite set \(S\subset\Z\), an \emph{\(S\)-packing center set} is
a set \(A\subseteq\Z\) such that the translates \(x+S\), \(x\in A\),
are pairwise disjoint. Put
\begin{equation}\label{eq:density}
 \ud(A)=\limsup_{N\to\infty}
 \frac{|A\cap[-N,N]|}{2N+1},\qquad
 \packd(S)=\sup\{\ud(A):A\text{ is an }S\text{-packing center set}\}.
\end{equation}
Here \(N\) ranges over the nonnegative integers. Thus density counts
centers, not covered integers. The definition permits nonperiodic sets
and does not require a density limit to exist.

\begin{theorem}\label{thm:main}
Let \(a,b\) be coprime positive integers, and let
\[
 S=\{0,a,2a,2a+b,2a+2b\}.
\]
Write the Euclidean division of \(a-b\) by \(5\) as
\[
 a-b=5q+r,\qquad q\in\Z,\quad 0\le r\le4.
\]
Then
\begin{equation}\label{eq:formula}
 \packd(S)=R(a,b):=
 \begin{cases}
 \dfrac15,&r=0,\\[4pt]
 \dfrac{b+2q}{2a+3b},&r=1,2,\\[6pt]
 \dfrac{a-2(q+1)}{3a+2b},&r=3,4.
 \end{cases}
\end{equation}
Every \(S\)-packing center set has symmetric upper density at most
\(R(a,b)\), and a periodic center set has density limit \(R(a,b)\).
\end{theorem}

Natalchenko and Sagdeev~\cite{NS} consider sets with \(k\) successive
gaps of size \(a\), followed by \(m\) gaps of size \(b\). Their
Theorem~2 supplies a lower bound for all positive \(k,m\), and the
unnumbered conjecture immediately following that theorem asserts its
sharpness. Equation~\eqref{eq:formula} is exactly the specialization
\(k=m=2\). Their Theorem~3 proves sharpness when \(k=1\) or \(m=1\).
The present theorem concerns only \(k=m=2\); it does not settle the
general conjecture. It is also distinct from their numbered
Conjecture~1, whose packing set has four elements.

The lower bounds and the equality case \(r=0\) already occur in
\cite[Theorem~2 and Section~2]{NS}. We include their specialized
constructions to make the proof self-contained. The upper bound below
uses an explicit path identity and its equality condition. No claim of
publication priority is made.

\section{Forbidden differences, symmetry, and the zero residue}

Define the finite set of positive integers
\begin{equation}\label{eq:M}
 M=\{ia+jb:0\le i,j\le2,\ i+j>0\}.
\end{equation}
A set \(A\subseteq\Z\) is \emph{\(M\)-avoiding} if
\(|x-y|\notin M\) whenever \(x,y\in A\) are distinct.
All coefficient indices in~\eqref{eq:M} are integers; repeated
numerical values are included only once.

The positive differences of elements of \(S\) are exactly
\[
 \{a,2a,b,2b,a+b,a+2b,2a+b,2a+2b\}=M.
\]
Indeed, the differences from \(0\) give
\(a,2a,2a+b,2a+2b\); those from \(a\) give
\(a,a+b,a+2b\); those from \(2a\) give \(b,2b\); and the
remaining difference is \(b\). Two distinct translates \(x+S\) and
\(y+S\) meet precisely when \(x-y=s'-s\) for some distinct
\(s,s'\in S\). Consequently, being an \(S\)-packing center set is
equivalent to being \(M\)-avoiding.

Exchanging \(a\) and \(b\) leaves \(M\) unchanged. For \(r\ne0\),
the exchanged quotient and remainder are
\begin{equation}\label{eq:swap}
 q'=-q-1,\qquad r'=5-r.
\end{equation}
If \(r=1,2\), the exchanged third branch equals
\[
 \frac{b-2(q'+1)}{3b+2a}=\frac{b+2q}{2a+3b}.
\]
If \(r=3,4\), the exchanged second branch equals
\[
 \frac{a+2q'}{2b+3a}=\frac{a-2(q+1)}{3a+2b}.
\]
The zero branch is unchanged as well. Thus the target formula and the
class of admissible center sets are both invariant under this exchange.

The five elements of \(S\) are strictly increasing. If \(A\) is an
\(S\)-packing center set, all translates with centers in \([-N,N]\)
lie in \([-N,N+2a+2b]\). Their disjointness gives
\begin{equation}\label{eq:volume}
 5|A\cap[-N,N]|\le2N+2a+2b+1.
\end{equation}
After division by \(2N+1\) and passage to the upper limit, this proves
\(\ud(A)\le1/5\).

If \(r=0\), then \(a\equiv b\equiv v\pmod5\) for some nonzero
residue \(v\), since otherwise \(5\) would divide both \(a\) and
\(b\). The residues of \(S\) are \(0,v,2v,3v,4v\), all distinct.
It follows that \(5\Z\) is an \(S\)-packing center set: an equality
\(5x+s=5y+s'\) first forces \(s=s'\) by residues and then forces
\(x=y\). Its density is \(1/5\), proving the zero branch. The only
coprime case with \(a=b\) is \(a=b=1\), which is included here.

For the upper bounds in the nonzero branches, symmetry now permits us
to assume \(a>b>0\). The next two lemmas do not require coprimality.

\section{A path identity and its equality condition}

Throughout this section, let \(a>b>0\), let \(A\) be \(M\)-avoiding,
and set
\begin{equation}\label{eq:windows}
 h=a+b,\quad n_1=2h+b=2a+3b,\quad n_2=2h+a=3a+2b,
 \qquad F_i=|A\cap[0,n_i)|\ (i=1,2).
\end{equation}
In particular, \(0<n_1<n_2\). Let \(W\) be the six words with two
letters \(a\) and two letters \(b\). For
\(w=(d_1,d_2,d_3,d_4)\in W\), define
\[
 p_0(w)=0,\qquad p_j(w)=\sum_{i=1}^j d_i\quad(1\le j\le4),
 \qquad P_w=\{p_0(w),\ldots,p_4(w)\}.
\]
The partial sums are strictly increasing. The positive difference of
any two is a sum of a nonempty consecutive part of the word, and hence
belongs to \(M\). Therefore
\begin{equation}\label{eq:path-clique}
 |A\cap(t+P_w)|\le1\qquad(t\in\Z,\ w\in W).
\end{equation}
For a subset \(E\) of \(\Z\), write \(\ind_E\) for its indicator
function on \(\Z\).

\begin{lemma}[Path identity and local bound]\label{lem:path}
The following identity holds pointwise on \(\Z\):
\begin{equation}\label{eq:path-identity}
 \sum_{w\in W}\sum_{t=0}^{d_4-1}\ind_{t+P_w}
 =3\ind_{[0,n_1)}+3\ind_{[0,n_2)}.
\end{equation}
Consequently,
\begin{equation}\label{eq:sum-bound}
 F_1+F_2\le h.
\end{equation}
If equality holds in~\eqref{eq:sum-bound}, then
\(|A\cap(t+P_w)|=1\) for every \(w\in W\) and every
\(0\le t<d_4\).
\end{lemma}

\begin{proof}
Since the partial sums are distinct, expanding the indicator of each
path and summing over its translations gives
\[
 \sum_{w\in W}\sum_{t=0}^{d_4-1}\ind_{t+P_w}
 =\sum_{j=0}^4\sum_{w\in W}\ind_{[p_j(w),p_j(w)+d_4)}.
\]
Fix \(j\in\{0,1,2,3\}\). Interchanging positions \(j+1\) and
\(4\) is a bijection of \(W\). It leaves the first \(j\) letters
unchanged and hence preserves \(p_j(w)\), while interchanging
\(d_{j+1}\) and \(d_4\). Thus
\[
 \sum_{w\in W}\ind_{[p_j(w),p_j(w)+d_4)}
 =\sum_{w\in W}\ind_{[p_j(w),p_j(w)+d_{j+1})}.
\]
For each fixed word, the four intervals on the right, as \(j\) varies
from \(0\) to \(3\), partition \([0,2h)\). Their total contribution
over the six words is \(6\ind_{[0,2h)}\). For \(j=4\), all partial
sums equal \(2h\); three words end with \(a\) and three end with
\(b\). The remaining contribution is
\[
 3\ind_{[2h,2h+a)}+3\ind_{[2h,2h+b)}.
\]
Adding these contributions proves~\eqref{eq:path-identity}.

Sum the identity over points of \(A\). All functions involved have
finite support, so this is a finite sum even when \(A\) is infinite.
There are \(3a+3b=3h\) indexed translated paths on the left. By
\eqref{eq:path-clique}, each contributes at most one, whence
\(3(F_1+F_2)\le3h\). If equality holds, these \(3h\) quantities,
each either zero or one, must all be one. Possible overlaps between
different indexed paths do not affect this argument.
\end{proof}

\begin{lemma}[Parity at equality]\label{lem:parity}
If \(F_1+F_2=h\), then there is a nonnegative integer \(e\) such that
\begin{equation}\label{eq:parity}
 F_1=b+2e,\qquad F_2=a-2e.
\end{equation}
\end{lemma}

\begin{proof}
Put \(x_z=\ind_A(z)\). The words \(abba\) and \(baba\) both end
with \(a\). Lemma~\ref{lem:path} therefore shows that, for every
\(0\le u<a\), each of the two sets
\[
 u+\{0,a,h,a+2b,2h\},\qquad
 u+\{0,b,h,a+2b,2h\}
\]
contains exactly one point of \(A\). Subtracting their indicator sums
yields \(x_{u+a}=x_{u+b}\). Summing over \(0\le u<a\), we obtain
\[
 |A\cap[a,2a)|=|A\cap[b,h)|.
\]
The endpoint order is \(b<a<h<2a\). Removing the common interval
\([a,h)\) from this equality gives
\begin{equation}\label{eq:equal-gaps}
 |A\cap[b,a)|=|A\cap[h,2a)|=:e.
\end{equation}

The word \(aabb\) has path \(S\) and ends with \(b\). Hence every
\(u+S\), \(0\le u<b\), contains exactly one point of \(A\).
These \(b\) sets are pairwise disjoint: their five point layers, as
\(u\) varies, are the disjoint intervals
\[
 [0,b),\quad [a,h),\quad [2a,2a+b),\quad
 [2a+b,2a+2b),\quad [2a+2b,2a+3b).
\]
Their union is
\([0,b)\mathbin{\dot\cup}[a,h)\mathbin{\dot\cup}[2a,n_1)\),
and therefore contains exactly \(b\) points of \(A\). Its complement
within \([0,n_1)\) consists of \([b,a)\) and \([h,2a)\). By
\eqref{eq:equal-gaps}, this proves \(F_1=b+2e\). The equality
\(F_1+F_2=a+b\) then gives \(F_2=a-2e\).
\end{proof}

Neither lemma assumes \(0\in A\). This will allow their use at every
integer starting point.

\section{The four nonzero residues and the global upper bound}

Keep \(a>b>0\), and write \(a-b=5q+r\) with \(1\le r\le4\).
Here \(q\ge0\).

\begin{proposition}\label{prop:local}
For every integer \(x\), at least one length
\(\ell\in\{n_1,n_2\}\) satisfies
\begin{equation}\label{eq:good-window}
 |A\cap[x,x+\ell)|\le R(a,b)\ell.
\end{equation}
\end{proposition}

\begin{proof}
First take \(x=0\).

If \(r=1\) or \(r=2\), we claim that
\begin{equation}\label{eq:alternative12}
 F_1\le b+2q\quad\text{or}\quad F_2\le b+3q.
\end{equation}
If both inequalities fail, integer-valuedness gives
\[
 F_1+F_2\ge(b+2q+1)+(b+3q+1)=h+2-r.
\]
For \(r=1\), this contradicts Lemma~\ref{lem:path}. For \(r=2\),
Lemma~\ref{lem:path} forces equality throughout. In particular,
\(F_1=b+2q+1\), whereas Lemma~\ref{lem:parity} requires
\(F_1\equiv b\pmod2\). This is a contradiction, proving
\eqref{eq:alternative12}. Set \(\rho=(b+2q)/n_1=R(a,b)\).
Substitution of \(a=b+5q+r\) gives
\begin{equation}\label{eq:cross12}
 (b+2q)n_2-(b+3q)n_1=br>0.
\end{equation}
Thus \((b+3q)/n_2<\rho\), so either alternative
in~\eqref{eq:alternative12} proves~\eqref{eq:good-window}.

If \(r=3\) or \(r=4\), the required integer alternative is
\begin{equation}\label{eq:alternative34}
 F_1\le a-3(q+1)\quad\text{or}\quad F_2\le a-2(q+1).
\end{equation}
Failure of both would give
\[
 F_1+F_2\ge\bigl(a-3(q+1)+1\bigr)
                +\bigl(a-2(q+1)+1\bigr)=h+r-3.
\]
For \(r=4\), this exceeds the bound in Lemma~\ref{lem:path}.
For \(r=3\), equality is forced, and
\[
 F_1=a-3(q+1)+1=b+2q+1,
\]
contrary to Lemma~\ref{lem:parity}. This proves
\eqref{eq:alternative34}. Now put
\(\rho=(a-2(q+1))/n_2=R(a,b)\). Direct expansion gives
\begin{equation}\label{eq:cross34}
 \bigl(a-2(q+1)\bigr)n_1-
 \bigl(a-3(q+1)\bigr)n_2=a(5-r)>0.
\end{equation}
It follows that \((a-3(q+1))/n_1<\rho\). Again either alternative
proves~\eqref{eq:good-window} at the origin.

For general \(x\), apply the result just obtained to
\(A-x=\{z-x:z\in A\}\). Translation preserves all differences,
so \(A-x\) is \(M\)-avoiding. Its count in \([0,\ell)\) is
exactly the count of \(A\) in \([x,x+\ell)\).
\end{proof}

\begin{lemma}[Bounded-length interval decomposition]\label{lem:global}
Let \(A\subseteq\Z\), let \(1\le u\le v\) be integers, and let
\(\rho\ge0\). Suppose that for every integer \(x\), some
\(\ell\in\{u,v\}\) satisfies
\(|A\cap[x,x+\ell)|\le\rho\ell\).
Then, for every nonnegative integer \(N\),
\begin{equation}\label{eq:finite-global}
 |A\cap[-N,N]|\le\rho(2N+1)+v.
\end{equation}
Consequently \(\ud(A)\le\rho\).
\end{lemma}

\begin{proof}
Start at the left endpoint of \([-N,N+1)\). Whenever its unprocessed
part has length at least \(v\), choose a good interval supplied by the
hypothesis and continue at its right endpoint. Its length is positive
and at most \(v\), so it remains within the unprocessed part and the
procedure terminates. The selected intervals are disjoint. The final
remainder has an integer length \(s\) with \(0\le s<v\). This
description also applies when \(2N+1<v\), in which case no interval
is selected.

The selected intervals contain at most \(\rho(2N+1-s)\) points of
\(A\); the remainder contains at most \(s\). Thus
\[
 |A\cap[-N,N]|\le\rho(2N+1-s)+s
 =\rho(2N+1)+(1-\rho)s\le\rho(2N+1)+v,
\]
using \(\rho\ge0\) and \(s<v\). Divide by \(2N+1\) and take
the upper limit. The constant \(v\) is fixed, and
\(v/(2N+1)\to0\), proving the density bound.
\end{proof}

To apply the lemma, note that for \(r=1,2\) the numerator \(b+2q\)
is positive; for \(r=3,4\),
\(a-2(q+1)=b+3q+r-2>0\). Proposition~\ref{prop:local} and
Lemma~\ref{lem:global}, with \(u=n_1\), \(v=n_2\), and
\(\rho=R(a,b)\), therefore give
\begin{equation}\label{eq:upper-all}
 \ud(A)\le R(a,b).
\end{equation}
The argument holds for every \(M\)-avoiding set, including infinite
nonperiodic sets, finite sets, and the empty set. Its error constant
depends only on the fixed parameters. Taking the supremum over center
sets proves the desired upper bound. Symmetry from
Section~2 supplies the case \(a<b\), and the case \(r=0\) was
already proved there.

\section{Periodic constructions attaining the bound}

We now verify the \(k=m=2\) specialization of the construction
in~\cite[Section~2]{NS}. In this section \(a,b\) are again any
coprime positive integers; no order is imposed on them, so the
quotient \(q\) may be negative.

Suppose first that \(r\in\{1,2\}\), and put
\begin{equation}\label{eq:construction-data}
 g=a-b,\qquad n=2a+3b,\qquad
 c=b+2q=\frac{n-2r}{5}.
\end{equation}
The integer \(c\) is positive because \(n\ge5>2r\). Define
\begin{equation}\label{eq:construction}
 A_0=\{tg:0\le t<c\}+n\Z,
\end{equation}
where \(t\) is an integer. We will prove that the indicated residues
are distinct modulo \(n\), and that \(A_0\) is \(M\)-avoiding.

Consider an integer relation
\begin{equation}\label{eq:relation}
 tg=ia+jb+zn,\qquad |t|<c,\quad 0\le i,j\le2,\quad z\in\Z.
\end{equation}
Rearranging it yields
\[
 a(t-2z-i)=b(t+3z+j).
\]
Since \(\gcd(a,b)=1\), Euclid's divisibility lemma implies that
\(a\) divides \(t+3z+j\). Writing this latter integer as \(ap\)
and cancelling the positive factor \(a\), we obtain an integer
\(p\) such that
\[
 t-2z-i=bp,\qquad t+3z+j=ap.
\]
Let \(u=z-qp\). Subtracting these equations and substituting
\(a-b=5q+r\) gives the first equality below; the first equation
above and \(c=b+2q\) give the second:
\begin{equation}\label{eq:division-relation}
 rp=5u+i+j,\qquad t=cp+2u+i.
\end{equation}
As \(0\le i+j\le4\), the first equality is Euclidean division by
\(5\). Thus \(u=\lfloor rp/5\rfloor\), also when \(p<0\).

If \(p>0\), then \(u\ge0\), so the second equality gives
\(t\ge cp\ge c\), contradicting \(|t|<c\).
If \(p<0\), then \(u\le-1\). Since \(i\le2\), we have
\(2u+i\le0\), and hence \(t\le cp\le-c\), another
contradiction. If \(p=0\), the first equality, together with
\(0\le i+j\le4\), forces \(u=0\) and \(i=j=0\). The definitions
then give \(z=0\) and \(t=0\). Therefore~\eqref{eq:relation} has
only the trivial solution \(t=i=j=z=0\).

If two of the residues in~\eqref{eq:construction} were equal modulo
\(n\), the difference of their indices would supply a relation
\eqref{eq:relation} with \(i=j=0\) and a nonzero \(t\) satisfying
\(|t|<c\). This is impossible. If two points of \(A_0\) had
positive difference \(ia+jb\in M\), their indices and period
shifts would give~\eqref{eq:relation} with \(i+j>0\), also
impossible. Negative forbidden differences are covered by reversing
the two points. This proves avoidance for the whole periodic set,
including points in different periods.

There are exactly \(c\) occupied residue classes modulo \(n\).
For any one residue class, its count in an interval of \(L\)
consecutive integers differs from \(L/n\) by less than one. Summing
over the \(c\) classes gives
\[
 \left||A_0\cap[-N,N]|-\frac{c}{n}(2N+1)\right|\le c.
\]
Its symmetric density limit is therefore \(c/n\), which is the
second branch of~\eqref{eq:formula}. The argument has used no
restriction on the sign of \(q\).

If \(r\in\{3,4\}\), exchange \(a\) and \(b\). The new remainder
\(r'=5-r\) belongs to \(\{1,2\}\), and the new quotient is
\(q'=-q-1\). Applying the construction just proved gives an
\(M\)-avoiding set of period \(2b+3a\) with density
\[
 \frac{a+2q'}{2b+3a}=
 \frac{a-2(q+1)}{3a+2b}.
\]
The forbidden difference set is unchanged by the exchange, so this is
a valid packing center set for the original \(S\). The zero branch
is attained by \(5\Z\) as proved in Section~2. All lower bounds
match~\eqref{eq:upper-all} and~\eqref{eq:volume}, completing the
proof of Theorem~\ref{thm:main}.

\section*{Computation and AI assistance}

The proof is the finite combinatorial and arithmetic argument given
above, followed by the explicit upper-limit estimate. It does not
depend on checking finitely many parameter pairs or on a theorem
asserting the existence of an optimal periodic set. Supporting exact
computations distributed with this note check only their stated
finite ranges. They are supplementary verification, not the argument
for the unrestricted parameters. This note does not assert that a
complete proof-assistant formalization has been obtained.

This work was developed with AI assistance, including assistance with
the mathematical argument, exposition, and supporting verification
code.

\begin{thebibliography}{9}
\bibitem{NS}
Alexander Natalchenko and Arsenii Sagdeev,
\emph{Packing Density of Sets With Only Two Nonmixed Gaps},
Journal of Graph Theory \textbf{110} (2025), no.~3, 263--271.
\href{https://doi.org/10.1002/jgt.23269}{doi:10.1002/jgt.23269}.
Version consulted: \href{https://arxiv.org/abs/2407.01101v2}
{arXiv:2407.01101v2}, revised June~6, 2025.
\end{thebibliography}

\end{document}
